\chapter{Higher-order spatial derivatives}
\label{chap:higher-order-spatial-derivatives}

\begin{quote}
Spec dir:
\passthrough{\lstinline!specs/2026-05-02-phase-8-higher-order-derivatives/!}.
Version: \passthrough{\lstinline!0.8.0!}.
\end{quote}

\subsection{What this gives you}\label{what-this-gives-you}

Every distinct scalar partial derivative on the periodic Cartesian grid
up to total rank 4, plus the biharmonic operator \(\nabla^4\). Each
operator is templated on accuracy \passthrough{\lstinline!Order = 2!}
(default) or \passthrough{\lstinline!Order = 4!}, mirroring Phase 7's
switching mechanism on \passthrough{\lstinline!laplacian!},
\passthrough{\lstinline!gradient!},
\passthrough{\lstinline!divergence!}.

\begin{lstlisting}[language={C++}]
diff::dxx(psi);              // ∂²ψ/∂x² (rank 2, order 2 default)
diff::dxy<4>(psi);           // ∂²ψ/∂x∂y (rank 2, order 4)
diff::d3dx2dy(psi);          // ∂³ψ/∂x²∂y (rank 3)
diff::d3dxdydz<4>(psi);      // ∂³ψ/∂x∂y∂z (rank 3, order 4)
diff::d4dx2dy2(psi);         // ∂⁴ψ/∂x²∂y² (rank 4)
diff::d4dx2dydz<4>(psi);     // ∂⁴ψ/∂x²∂y∂z (rank 4, order 4)
diff::biharmonic(psi);       // ∇⁴ψ (rank 4, contracted)
diff::d4(psi);               // alias for biharmonic — same result
\end{lstlisting}

\subsection{Naming convention}\label{naming-convention}

\passthrough{\lstinline!d<rank>d<axis-with-exponent>...!}. The rules:

\begin{enumerate}
\def\labelenumi{\arabic{enumi}.}
\tightlist
\item
  \textbf{Total rank} appears once after the leading
  \passthrough{\lstinline!d!} (e.g., \passthrough{\lstinline!d3!} for a
  rank-3 partial).
\item
  \textbf{Each distinct axis} is listed in \textbf{lexicographic order}
  (\passthrough{\lstinline!x!} → \passthrough{\lstinline!y!} →
  \passthrough{\lstinline!z!}), with its \textbf{exponent} (the
  multiplicity of differentiation along that axis). Exponent of
  \passthrough{\lstinline!1!} is omitted.
\item
  Concatenated as one identifier; no separators, lowercase.
\end{enumerate}

{\def\LTcaptype{none} % do not increment counter
\begin{longtable}[]{@{}
  >{\raggedright\arraybackslash}p{(\linewidth - 2\tabcolsep) * \real{0.7541}}
  >{\raggedright\arraybackslash}p{(\linewidth - 2\tabcolsep) * \real{0.2459}}@{}}
\toprule\noalign{}
\begin{minipage}[b]{\linewidth}\raggedright
Math notation
\end{minipage} & \begin{minipage}[b]{\linewidth}\raggedright
Function name
\end{minipage} \\
\midrule\noalign{}
\endhead
\bottomrule\noalign{}
\endlastfoot
\(\partial^2 \psi / \partial x^2\) & \passthrough{\lstinline!dxx!} \\
\(\partial^2 \psi / \partial y^2\) & \passthrough{\lstinline!dyy!} \\
\(\partial^2 \psi / \partial z^2\) & \passthrough{\lstinline!dzz!} \\
\(\partial^2 \psi / \partial x \partial y\) &
\passthrough{\lstinline!dxy!} \\
\(\partial^2 \psi / \partial x \partial z\) &
\passthrough{\lstinline!dxz!} \\
\(\partial^2 \psi / \partial y \partial z\) &
\passthrough{\lstinline!dyz!} \\
\(\partial^3 \psi / \partial x^3\) & \passthrough{\lstinline!d3dx3!} \\
\(\partial^3 \psi / \partial x^2 \partial y\) &
\passthrough{\lstinline!d3dx2dy!} \\
\(\partial^3 \psi / \partial x \partial y^2\) &
\passthrough{\lstinline!d3dxdy2!} \\
\(\partial^3 \psi / \partial x \partial y \partial z\) &
\passthrough{\lstinline!d3dxdydz!} \\
\(\partial^4 \psi / \partial x^4\) & \passthrough{\lstinline!d4dx4!} \\
\(\partial^4 \psi / \partial x^3 \partial y\) &
\passthrough{\lstinline!d4dx3dy!} \\
\(\partial^4 \psi / \partial x^2 \partial y^2\) &
\passthrough{\lstinline!d4dx2dy2!} \\
\(\partial^4 \psi / \partial x^2 \partial y \partial z\) &
\passthrough{\lstinline!d4dx2dydz!} \\
\(\nabla^4 \psi\) (contracted) & \passthrough{\lstinline!biharmonic!}
\emph{or} \passthrough{\lstinline!d4!} \\
\end{longtable}
}

A single shorthand like \passthrough{\lstinline!d4!} (no axis suffix)
denotes the \textbf{contracted scalar operator} at that rank ---
\(\nabla^4\) for \passthrough{\lstinline!d4!}. Phase 8 ships only the
even-rank shorthand \passthrough{\lstinline!d4!} because \$\nabla\^{}2 =
\$ Laplacian already has the familiar name
\passthrough{\lstinline!laplacian!} and isn't being renamed.

\subsection{API}\label{api}

Headers under
\href{../../../include/gridcalc/diff/}{\passthrough{\lstinline!include/gridcalc/diff/!}}:

{\def\LTcaptype{none} % do not increment counter
\begin{longtable}[]{@{}
  >{\raggedright\arraybackslash}p{(\linewidth - 2\tabcolsep) * \real{0.2243}}
  >{\raggedright\arraybackslash}p{(\linewidth - 2\tabcolsep) * \real{0.7757}}@{}}
\toprule\noalign{}
\begin{minipage}[b]{\linewidth}\raggedright
Header
\end{minipage} & \begin{minipage}[b]{\linewidth}\raggedright
Operators
\end{minipage} \\
\midrule\noalign{}
\endhead
\bottomrule\noalign{}
\endlastfoot
\passthrough{\lstinline!MixedPartial.hpp!} &
\passthrough{\lstinline!dxx!}, \passthrough{\lstinline!dyy!},
\passthrough{\lstinline!dzz!}, \passthrough{\lstinline!dxy!},
\passthrough{\lstinline!dxz!}, \passthrough{\lstinline!dyz!} (6 rank-2
partials) \\
\passthrough{\lstinline!ThirdOrder.hpp!} & 10 rank-3 partials
(\passthrough{\lstinline!d3dx3!}, \ldots,
\passthrough{\lstinline!d3dxdydz!}) \\
\passthrough{\lstinline!FourthOrder.hpp!} & 15 rank-4 partials
(\passthrough{\lstinline!d4dx4!}, \ldots,
\passthrough{\lstinline!d4dxdydz2!}) \\
\passthrough{\lstinline!Biharmonic.hpp!} &
\passthrough{\lstinline!biharmonic!}, \passthrough{\lstinline!d4!}
alias \\
\end{longtable}
}

Every function has the same signature shape:

\begin{lstlisting}[language={C++}]
template <int Order = 2>
inline core::Field<double> NAME(const core::Field<double>& psi);
\end{lstlisting}

\passthrough{\lstinline!Order!} must be a supported even integer
(\passthrough{\lstinline!2!} or \passthrough{\lstinline!4!} at Phase 8).
Calling e.g.~\passthrough{\lstinline!d3dx2dy<3>(psi)!} triggers a
\textbf{compile error} at the call site (the underlying primary template
is intentionally undefined), not a silent zero-stencil bug.

\subsubsection{Stencil radius and per-point
cost}\label{stencil-radius-and-per-point-cost}

{\def\LTcaptype{none} % do not increment counter
\begin{longtable}[]{@{}
  >{\raggedright\arraybackslash}p{(\linewidth - 6\tabcolsep) * \real{0.2321}}
  >{\raggedright\arraybackslash}p{(\linewidth - 6\tabcolsep) * \real{0.2054}}
  >{\raggedright\arraybackslash}p{(\linewidth - 6\tabcolsep) * \real{0.2054}}
  >{\raggedright\arraybackslash}p{(\linewidth - 6\tabcolsep) * \real{0.3571}}@{}}
\toprule\noalign{}
\begin{minipage}[b]{\linewidth}\raggedright
Operator family
\end{minipage} & \begin{minipage}[b]{\linewidth}\raggedright
Order 2 radius / axis
\end{minipage} & \begin{minipage}[b]{\linewidth}\raggedright
Order 4 radius / axis
\end{minipage} & \begin{minipage}[b]{\linewidth}\raggedright
Per-point reads (worst case at rank 4)
\end{minipage} \\
\midrule\noalign{}
\endhead
\bottomrule\noalign{}
\endlastfoot
Single-axis (\passthrough{\lstinline!d3dx3!},
\passthrough{\lstinline!d4dx4!}) & 2 / 0 / 0 & 3 / 0 / 0 &
\passthrough{\lstinline!5!} (order 2) → \passthrough{\lstinline!7!}
(order 4) \\
Two-axis (\passthrough{\lstinline!d3dx2dy!},
\passthrough{\lstinline!d4dx2dy2!}) & 2 / 1 / 0 → 2 / 2 / 0 & 3 / 2 / 0
→ 3 / 3 / 0 & up to \passthrough{\lstinline!25!} (order 2) →
\passthrough{\lstinline!49!} (order 4) \\
Three-axis (\passthrough{\lstinline!d3dxdydz!},
\passthrough{\lstinline!d4dx2dydz!}) & various & various & up to
\passthrough{\lstinline!27!} (order 2) → \passthrough{\lstinline!147!}
(order 4) \\
\passthrough{\lstinline!biharmonic!} & 6 stencils combined & 6 stencils
combined & order 2: \textasciitilde{}\passthrough{\lstinline!27!}; order
4: \textasciitilde{}\passthrough{\lstinline!81!} \\
\end{longtable}
}

Per-axis \(1/h_a^{N_a}\) scaling means \textbf{anisotropic per-axis
spacing} is handled by the existing \passthrough{\lstinline!Field!}'s
\passthrough{\lstinline!Vec3d \{hx, hy, hz\}!} without any weight-table
modification.

\subsection{\texorpdfstring{Worked example: \(\nabla^4\) on a sinusoidal
eigenfunction}{Worked example: \textbackslash nabla\^{}4 on a sinusoidal eigenfunction}}\label{worked-example-nabla4-on-a-sinusoidal-eigenfunction}

For \(\psi(\mathbf{x}) = \sin(\mathbf{k}\cdot\mathbf{x})\) on the
periodic \([0, 2\pi)^3\) box, the biharmonic recovers
\(|\mathbf{k}|^4\):

\[\nabla^4 \psi = |\mathbf{k}|^4\,\psi.\]

Numerically:

\begin{lstlisting}[language={C++}]
#include <cmath>
#include <gridcalc/core/EigenAliases.hpp>
#include <gridcalc/core/Field.hpp>
#include <gridcalc/core/Grid.hpp>
#include <gridcalc/diff/Biharmonic.hpp>

using gridcalc::core::Field;
using gridcalc::core::Grid;
using gridcalc::core::Vec3d;
using gridcalc::diff::biharmonic;

constexpr double kPi = 3.14159265358979323846;

const int N = 64;
const double L = 2.0 * kPi;
const double h = L / N;

Grid grid(N, N, N, Vec3d(h, h, h));
Field<double> psi(grid, [](double x, double y, double z) {
  return std::sin(x + 2.0 * y + z);                    // k = (1, 2, 1), |k|² = 6
});

Field<double> bh_o2 = biharmonic<2>(psi);              // ~6% relative error at N=64
Field<double> bh_o4 = biharmonic<4>(psi);              // ~3e-3 relative error at N=64

// bh ≈ |k|⁴ * psi = 36 * psi
\end{lstlisting}

For the same accuracy, order 4 lets you stay on a coarser grid and saves
substantial memory in 3D (each grid refinement is 8× the storage).

\subsection{Combining with Phase 1--7
operators}\label{combining-with-phase-17-operators}

The Phase 8 operators are first-class members of
\passthrough{\lstinline!gridcalc::diff!} and compose with everything
from Phases 1--7 by ordinary function composition. Examples:

\begin{itemize}
\tightlist
\item
  A \textbf{biharmonic regularizer} on a free-energy functional written
  through Phase 4's \passthrough{\lstinline!func::evaluate!}: pass
  \passthrough{\lstinline![](const Field\& f) \{ return biharmonic<4>(f); \}!}
  as the integrand.
\item
  A \textbf{fourth-order PDE driver} through Phase 6's
  \passthrough{\lstinline!solve::integrate!}: build the RHS callable
  using \passthrough{\lstinline!diff::biharmonic<4>!} plus any other
  Phase 8 partials.
\end{itemize}

\subsection{\texorpdfstring{What this does \emph{not} do
(yet)}{What this does not do (yet)}}\label{what-this-does-not-do-yet}

\begin{itemize}
\tightlist
\item
  \textbf{No order ≥ 6 stencils.} Same call as Phase 7 --- needs a
  Fornberg generator the project has not yet built.
\item
  \textbf{No order-3 / asymmetric stencils.} Symmetric central
  differences have even truncation orders only.
\item
  \textbf{No vector- or tensor-valued partial-derivative entry points.}
  The Phase 8 family is
  \passthrough{\lstinline!Field<double> → Field<double>!} only.
  Up-to-4th-order \emph{gradient} (Phase 11) and vector / tensor fields
  (Phase 13/14) own those.
\item
  \textbf{No spectral verification.} Phase 9 adds the FFT cross-check on
  every Phase 8 operator.
\item
  \textbf{No non-orthogonal Bravais grids.} Phase 15 ---
  cross-derivative terms appear and the per-axis-independent stencil
  structure breaks down.
\item
  \textbf{No runtime-order or runtime-multi-index dispatch.}
  Compile-time only; add a thin wrapper layer downstream if
  config-file-driven dispatch is needed.
\end{itemize}

For the Taylor-matching derivation of the 3rd- and 4th-derivative weight
tables, the outer-product reasoning behind the multi-axis stencils, and
the design rationale for the d-prefix naming convention, see
\href{../../developer-note/notes/phase-8-higher-order-derivatives.md}{\passthrough{\lstinline!docs/developer-note/notes/phase-8-higher-order-derivatives.md!}}.
